% Quarantine-only source-aware proof.
% Date: 29 August 2026.
% This file is not a live-book input and makes no novelty or priority claim.
% DEPENDENCY BLOCK: the proof is conditional on a corrected and verified
% arbitrary normal-curve classification.

\begin{proof}[Proof of
\cref{prop:p213-componentwise-arbitrary-valuation}]
We give the reductions in full because the nilpotent calculus, generic
multiplicities, classification, and curvature measures remember different
parts of the analytic structure.

\paragraph{Nilreduction and the ambient geometry.}
An analytic seminorm kills every nilpotent.  Hence $X$ and
$X_{\mathrm{red}}$ have the same underlying topological space, the same
analytic points, and the same completed residue fields.  Affinoid charts of
the reduction are quotients of affinoid charts of $X$ by their nilradicals.
It follows directly that goodness, separatedness, strictness, and pure
dimension one pass to $X_{\mathrm{red}}$.  The immersion $\iota$ is finite
and has empty relative boundary by \cite[\S11.4.4]{CLD25}.  The
boundary-composition identity \cite[(11.4.1.3)]{CLD25} therefore shows that
$X_{\mathrm{red}}$ is boundaryless.

Both spaces are good and topologically separated, hence locally
holomorphically separated by \cite[\S15.2.2]{CLD25} and locally
differentially separated by \cite[Corollary~15.2.4]{CLD25}.  Thus the Radon
measure and Borel-multiplication statements below apply by
\cite[Lemma~7.2.4, Proposition~7.2.8, and Example~7.2.12]{CLD25}.

If $A\to A/I$ is an affinoid quotient with $I$ nilpotent, every unit of
$A/I$ lifts to a unit of $A$: if $ab=1+n$ with $n$ nilpotent, the finite
geometric series for $(1+n)^{-1}$ corrects $b$ to an inverse.  Every
analytic seminorm on $A$ kills $I$, so a lifted unit and its reduction have
the same absolute value at every analytic point.  Thus a moment map on a
reduced chart lifts, after shrinking, to a moment map on $X$ with the same
tropicalization, and every moment on $X$ restricts to one on the reduction.
This is the nilpotent-immersion case of
\cite[Proposition~15.1.3(b) and the following form variant]{CLD25}.

The construction of ordinary CLD forms from moment charts gives
\begin{equation}
 \label{eq:p213-componentwise-av-form-isomorphism}
 \iota^*:A_X^{p,q}\xrightarrow{\ \sim\ }
                    A_{X_{\mathrm{red}}}^{p,q}.
\end{equation}
The equality of tropicalizations makes
\eqref{eq:p213-componentwise-av-form-isomorphism} compatible with wedge
product, $J$, $\mathrm d'$, $\mathrm d''$, support, and every positivity
cone.  For the topology, fix a compact test support.  A compact-test datum
is a finite collection of moment charts and Lagerberg representatives.
Restriction and the lifted charts above are cofinal collections of such
data, and on corresponding charts the tropical images, subdivisions,
representatives, and all $C^r$-seminorms coincide.  Hence
\begin{equation}
 \label{eq:p213-componentwise-av-test-homeomorphism}
 \iota^*:A^{p,q}_{X,c}\xrightarrow{\ \sim\ }
                    A^{p,q}_{X_{\mathrm{red}},c}
\end{equation}
is a homeomorphism of locally convex spaces.

Taking continuous duals gives
\begin{equation}
 \label{eq:p213-componentwise-av-current-isomorphism}
 \iota_*:D^{a,b}(X_{\mathrm{red}})
       \xrightarrow{\ \sim\ }D^{a,b}(X),
 \qquad
 \langle\iota_*S,\alpha\rangle
       =\langle S,\iota^*\alpha\rangle.
\end{equation}
The test-form positivity cones agree, so this isomorphism preserves
positivity.  It preserves $\mathrm d'$- and $\mathrm d''$-closedness because
$\iota^*$ commutes with the two differentials.  It also identifies the weak
topologies and multiplication by smooth functions.

If $S$ is positive, its product with a locally bounded Borel function is
defined through the Radon measures attached to positive test forms, as in
\cite[Example~7.2.12]{CLD25}.  Those test forms and measures agree under
\eqref{eq:p213-componentwise-av-form-isomorphism}, so
$\iota_*(fS)=f\iota_*S$.  The same identity holds by linearity when locally
$S$ is a difference of two positive currents.  It is independent of the
decomposition: an equality between two sums of positive currents is an
equality of all attached measures, and remains so after multiplication by
$f$.  This does not define a Borel product for an arbitrary order-zero
current.

Finally, the first stage in the CLD psh filtration is the same on $X$ and
$X_{\mathrm{red}}$, because smoothness and positivity of
$\mathrm d'\mathrm d''u$ agree.  Successor stages take decreasing limits,
the continuous filtration takes local uniform limits, and limit stages
take local unions.  Transfinite induction identifies the psh and continuous
psh sheaves.  The $G$-form assertion follows from the $G$-form variant after
\cite[Proposition~15.1.3]{CLD25}.  This proves
\cref{itm:p213-componentwise-av:calculus}.

\paragraph{Normalization and the deleted locus.}
Let
\[
 \bar\nu:\coprod_i\widetilde X_i\longrightarrow X_{\mathrm{red}}
\]
be the componentwise normalization.  On a good affinoid domain define
\[
 \mathcal Q:=\bar\nu_*\mathcal O_{\coprod_i\widetilde X_i}
                  /\mathcal O_{X_{\mathrm{red}}},
 \qquad
 \mathfrak c:=\operatorname{Ann}_{\mathcal O_{X_{\mathrm{red}}}}
                  (\mathcal Q).
\]
Let $Z_0=\operatorname{Supp}(\mathcal Q)=V(\mathfrak c)$ and let $Z$ be
the union of $Z_0$ with the pairwise intersections of distinct irreducible
components.  Analytic local rings on a good analytic space are excellent by
\cite[\S11.1.10]{CLD25}; normalization is finite and its components
correspond to the irreducible components by \cite[\S11.3.1]{CLD25}.
Thus $\mathcal Q$ and the conductor are coherent.  Their support misses a
dense open subset of every reduced component and is zero-dimensional.
The same dimension argument applies to a proper intersection of two
components.  Irreducible components are locally finite by
\cite[\S11.2.5]{CLD25}; on a relatively compact affinoid neighborhood only
finitely many occur.  Hence $Z$ is closed and locally finite.

Put
\begin{equation}
 \label{eq:p213-componentwise-av-conductor-preimage}
                         Z_i:=\bar\nu_i^{-1}(Z).
\end{equation}
It is a proper zero-dimensional Zariski closed subset of the normal curve
$\widetilde X_i$, hence is locally finite.  We verify that it consists of
rigid points without using a nontrivial ground valuation or a compact
strict representative.  Suppose that $z\in Z_i$ were nonrigid.  For a
normal analytic curve, \cite[\S4.2.18]{Ducros24CurveStructure} gives
\[
                 \mathcal O_{\widetilde X_i,z}=\kappa(z),
\]
a field.  A coherent ideal defining $Z_i$ near $z$ has proper stalk in
this field and therefore has zero stalk.  Coherence then makes that ideal
zero on a neighborhood of $z$, forcing the supposedly zero-dimensional
closed set to contain a one-dimensional neighborhood, a contradiction.
Thus every point of $Z_i$ is rigid.  Every point of $Z$ has a preimage in
some $Z_i$.  A finite morphism induces a finite extension of completed
residue fields; since the upper field is finite over $K$, its subfield at
the image is finite over $K$ as well.  Thus every point of $Z$ is rigid.

The finite normalization remains strict in the sense used by CLD.  Indeed,
at a point above a strict germ, homogeneous degrees coming from the base
are torsion modulo $|K^\times|$ by
\cite[Lemma~11.8.5]{CLD25}.  Integrality for a finite local algebra makes
the degree of every additional homogeneous element torsion over those
base degrees; hence it too is torsion modulo $|K^\times|$.  The same lemma
then proves strictness of the normalization germ.  The maps
$X_i\to X_{\mathrm{red}}$ and $\widetilde X_i\to X_i$ are finite and have
empty relative boundary.  Formula (11.4.1.3) in \cite{CLD25} shows that
$\widetilde X_i$ is boundaryless.  It is good and separated, normal by
construction, and connected because it normalizes one irreducible
component.

Set
\begin{equation}
 \label{eq:p213-componentwise-av-normal-open}
                         U_i:=\widetilde X_i\setminus Z_i.
\end{equation}
This is nonempty because $Z_i$ is a proper closed subset.  It is connected:
at a rigid point of a normal analytic curve the normalization fibre has
cardinality one, so the point is unibranch by
\cite[Proposition~3.5.2]{Ducros24CurveStructure}.  A connected analytic
curve is connected by injective arcs by
\cite[Theorem~3.5.1 and \S1.3.5]{Ducros24CurveStructure}, while an interior
point of such an arc is pluribranch by
\cite[\S1.3.12]{Ducros24CurveStructure}.  An injective arc joining two
points outside $Z_i$ therefore cannot pass through a point of $Z_i$.
Thus $U_i$ is a nonempty connected good separated normal boundaryless
purely one-dimensional analytic space.  No smoothness property has been
used.

We shall also use that $\delta_{U_i}$ is nonzero.  This follows directly
from the support theorem for integration currents, without any variation
graph.  The space $U_i$ is boundaryless, topologically separated, and
purely one-dimensional, so \cite[Theorem~12.11.1]{CLD25} makes its
associated tropical space very rich.  Therefore
\cite[Corollary~7.3.13]{CLD25} gives
\begin{equation}
 \label{eq:p213-componentwise-av-nonzero-integration}
                         \operatorname{Supp}(\delta_{U_i})=U_i.
\end{equation}
Since $U_i$ is nonempty, $\delta_{U_i}\ne0$.

\paragraph{Generic multiplicities.}
By \cite[Definition~13.2.5 and \S13.5.2]{CLD25},
\[
                   [X]=\sum_i m_i[X_i],
 \qquad
                   [X_{\mathrm{red}}]=\sum_i[X_i].
\]
Integration agrees with the current attached to the fundamental cycle by
\cite[Proposition~16.6.3]{CLD25}.  The normalization is finite and has
calibrated degree one over the normal locus.  Its exceptional image is
contained in $Z$.  An ordinary smooth $(1,1)$-test germ vanishes at a rigid
point: by \cite[Proposition~15.1.1(b)]{CLD25} it is pulled back from a
local ordinary moment, and the dimension argument in the proof of
\cite[Proposition~15.2.1(a)]{CLD25} makes that pullback zero wherever
$d_K<1$.  A rigid point has $d_K=0$ by \cite[(11.2.3.1)]{CLD25}.
Consequently the cycle identities give the current identities
\eqref{eq:p213-componentwise-av-fundamental-cycles} on all of $X$.
Local finiteness of the components makes both sums locally finite.

If every $m_i$ is one, the two currents are equal.  Conversely, restrict an
asserted equality to $X_i\setminus Z$ and transport it to $U_i$.  It becomes
$(m_i-1)\delta_{U_i}=0$.  Equation
\eqref{eq:p213-componentwise-av-nonzero-integration} gives $m_i=1$.
Since $m_i$ is the length of the generic local ring along $X_i$, this is
equivalent to generic reducedness along every component.  This proves
\cref{itm:p213-componentwise-av:integration}.

\paragraph{Current classification and uniqueness.}
Let $T$ be a positive $\mathrm d'$- and $\mathrm d''$-closed current of
bidegree $(0,0)$ on $X$.  By
\eqref{eq:p213-componentwise-av-current-isomorphism}, there is a unique
positive closed current $T_{\mathrm{red}}$ on $X_{\mathrm{red}}$ with
$T=\iota_*T_{\mathrm{red}}$.  The open set
$X_{\mathrm{red}}\setminus Z$ is the disjoint union of the normal component
loci $X_i\setminus Z$, and
\[
 \bar\nu_i|_{U_i}:U_i\xrightarrow{\ \sim\ }X_i\setminus Z.
\]
Transport the restriction of $T_{\mathrm{red}}$ on $X_i\setminus Z$ to a
current $T_i$ on $U_i$.  Restriction preserves positivity and closedness by
\cite[\S7.1.6 and \S\S7.7.1--7.7.3]{CLD25}.  The arbitrary-valuation normal
curve classification
\cref{prop:p213-normal-curve-arbitrary-valuation-current} gives a number
$c_i\in\mathbf R_{\geq0}$ such that
\begin{equation}
 \label{eq:p213-componentwise-av-classification-off-Z}
                            T_i=c_i\delta_{U_i}.
\end{equation}
The coefficient is unique by
\eqref{eq:p213-componentwise-av-nonzero-integration}.

Write $\bar\delta_i=(\bar\nu_i)_*\delta_{\widetilde X_i}$ on
$X_{\mathrm{red}}$ and set
\[
                         S:=T_{\mathrm{red}}-\sum_i c_i\bar\delta_i.
\]
The sum is locally finite, and
\eqref{eq:p213-componentwise-av-classification-off-Z} says that $S$
vanishes on $X_{\mathrm{red}}\setminus Z$.  If $\alpha$ is a compactly
supported ordinary smooth $(1,1)$-test form, the ordinary moment and
dimension argument just used gives
\[
 \operatorname{Supp}(\alpha)
    \subseteq\{x:d_K(x)\geq1\}
    \subseteq X_{\mathrm{red}}\setminus Z.
\]
Thus $\langle S,\alpha\rangle=0$, and $S=0$.  Pushing forward by $\iota$
proves \eqref{eq:p213-componentwise-av-classification}.  Its uniqueness
follows after restriction to $X_i\setminus Z$, transport to $U_i$, and
testing against the form in
\eqref{eq:p213-componentwise-av-nonzero-integration}.  This proves
\cref{itm:p213-componentwise-av:classification}.

\paragraph{Componentwise curvature measures.}
Let $W,U,V$ be as in
\cref{itm:p213-componentwise-av:bedford-taylor}, and regard $W$ also as an
open subset of $X_{\mathrm{red}}$.  Put
\[
             \widetilde W_i:=\bar\nu_i^{-1}(W),
\]
and denote the restrictions of $\nu_i$ and $\bar\nu_i$ by
$\nu_{i,W}^X$ and $\nu_{i,W}^{\mathrm{red}}$, respectively.  The good
separated spaces $\widetilde X_i$ are locally differentially separated by
\cite[\S15.2.2 and Corollary~15.2.4]{CLD25}.  Pullback preserves continuous
psh functions by \cite[Proposition~8.1.5]{CLD25}.  Hence
\begin{equation}
 \label{eq:p213-componentwise-av-measures}
 \mu_i^U:=\mathrm d'\mathrm d''[\bar\nu_i^*U],
 \qquad
 \mu_i^V:=\mathrm d'\mathrm d''[\bar\nu_i^*V]
\end{equation}
are positive Radon measures by \cite[Lemma~8.5.5]{CLD25}.  Put
$\sigma_i:=\mu_i^U-\mu_i^V$.

For every continuous function $F$ on $W$, multiply the two identities in
\eqref{eq:p213-componentwise-av-fundamental-cycles} by $F$.  The product
is defined by \cite[Example~7.2.12]{CLD25}, and direct testing gives
\begin{align}
 [F]_X
 &=\sum_i m_i(\nu_{i,W}^X)_*
       \bigl(((\nu_{i,W}^X)^*F)\delta_{\widetilde W_i}\bigr),
 \label{eq:p213-componentwise-av-function-current-X}\\
 [F]_{X_{\mathrm{red}}}
 &=\sum_i(\nu_{i,W}^{\mathrm{red}})_*
       \bigl(((\nu_{i,W}^{\mathrm{red}})^*F)
                              \delta_{\widetilde W_i}\bigr).
 \label{eq:p213-componentwise-av-function-current-red}
\end{align}
Finite proper pushforward commutes with the current differentials by
transposition; for continuous psh functions this is also
\cite[Proposition~8.5.6]{CLD25}.  Differentiating
\eqref{eq:p213-componentwise-av-function-current-X} for $F=U,V$ shows
that the left equality in
\eqref{eq:p213-componentwise-av-curvature-equivalence} is equivalent to
\begin{equation}
 \label{eq:p213-componentwise-av-weighted-sum}
                 \sum_i m_i(\nu_{i,W}^X)_*\sigma_i=0.
\end{equation}
Since $\nu_{i,W}^X=\iota\circ\nu_{i,W}^{\mathrm{red}}$ and $\iota_*$ is
injective by \eqref{eq:p213-componentwise-av-current-isomorphism}, this is
equivalent to
\begin{equation}
 \label{eq:p213-componentwise-av-weighted-sum-red}
       \sum_i m_i(\nu_{i,W}^{\mathrm{red}})_*\sigma_i=0
       \qquad\text{on }X_{\mathrm{red}}.
\end{equation}
Likewise the right equality is equivalent to
\begin{equation}
 \label{eq:p213-componentwise-av-unweighted-sum}
                 \sum_i(\nu_{i,W}^{\mathrm{red}})_*\sigma_i=0.
\end{equation}

Set $Z_W=Z\cap W$ and $Z_{i,W}=Z_i\cap\widetilde W_i$.  Away from $Z_W$,
the component images are disjoint and each
$\nu_{i,W}^{\mathrm{red}}$ is an analytic isomorphism onto its component
image.  Restricting
\eqref{eq:p213-componentwise-av-weighted-sum-red} to one component image,
and using $m_i>0$, gives
\begin{equation}
 \label{eq:p213-componentwise-av-support}
                    \operatorname{Supp}(\sigma_i)\subseteq Z_{i,W}.
\end{equation}
The same conclusion follows from
\eqref{eq:p213-componentwise-av-unweighted-sum}.

Every point of $Z_{i,W}$ is rigid, and its local ring on
$\widetilde X_i$ is normal.  If $K$ is nontrivially valued, apply
\cref{lem:p21p3-componentwise-normal-rigid-zero-atom}; if $K$ is
trivially valued, apply
\cref{lem:p213-trivial-gauss-rigid-zero-atom}.  Applied to each of the two
continuous psh functions in
\eqref{eq:p213-componentwise-av-measures}, the appropriate lemma gives
\[
                         \mu_i^U(\{z\})
                         =\mu_i^V(\{z\})=0
                         \qquad(z\in Z_{i,W}).
\]
Because $Z_{i,W}$ is locally finite, a signed Radon measure supported there
is the locally finite sum of its point masses.  Equation
\eqref{eq:p213-componentwise-av-support} and the displayed zero-atom
equalities therefore imply
\begin{equation}
 \label{eq:p213-componentwise-av-component-equality}
                            \sigma_i=0
                            \qquad\text{for every }i.
\end{equation}
Starting from either the weighted or the unweighted sum yields
\eqref{eq:p213-componentwise-av-component-equality}; conversely that
equality implies both sums.  This proves
\eqref{eq:p213-componentwise-av-curvature-equivalence}.

Finally put $h=U-V$ and use
\eqref{eq:p213-componentwise-av-classification}.  Multiplication by $h$ is
defined because $T$ is positive and $h$ is locally bounded.  On compact
test supports the component sum is finite.  Proper pushforward,
\eqref{eq:p213-componentwise-av-component-equality}, and linearity give
\begin{align*}
 \mathrm d'\mathrm d''(hT|_W)
 &=\sum_i c_i(\nu_{i,W}^X)_*
       \mathrm d'\mathrm d''[\bar\nu_i^*U-\bar\nu_i^*V]\\
 &=\sum_i c_i(\nu_{i,W}^X)_*(\mu_i^U-\mu_i^V)=0.
\end{align*}
This is \eqref{eq:p213-componentwise-av-bt-conclusion} and proves
\cref{itm:p213-componentwise-av:bedford-taylor}.
\end{proof}
